\chapter{Физическая модель и редукция}\label{ch:model}

В главе строится модель обновления электрофоретического дисплея на трёх
уровнях абстракции: континуальном (уравнения Пуассона--Нернста--Планка),
усреднённом (система ОДУ для одного пикселя) и дискретном (готовом к
интеграции с задачей оптимального управления гл.~\ref{ch:theorem}). Каждый из
переходов мы сопровождаем допущениями, на которых он держится.
Глава завершается калибровкой эффективных параметров по данным
стенда (гл.~\ref{ch:experiment}).

\section{Микрокапсула и электрофорез}\label{sec:model-capsule}

Базовая структурная единица активной матрицы~--- \emph{микрокапсула}:
сферическая оболочка диаметром $d_{\mathrm{caps}}\sim 40\ldots 50$~мкм,
заполненная непроводящей жидкостью (углеводородная смесь с относительной
проницаемостью $\varepsilon_m \approx 2$). В ней диспергированы пигментные
частицы двух типов: белые (диоксид титана $\mathrm{TiO}_2$) и чёрные
(углеродная сажа) радиусом $R \sim 100\ldots 500$~нм
\cite{comiskey1998_nature, yang2021_mechanisms}. Заряд сообщают агенты
управления зарядом (\emph{charge controlling agents}, CCA)~--- ПАВ,
обменивающие заряды с поверхностью пигмента; при типичной концентрации CCA
заряд частицы~$q$ постоянен на временах одного
обновления~\cite{yang2021_mechanisms}.

Сверху и снизу к капсуле примыкают электроды: общий прозрачный из оксида
индия-олова и сегментированный нижний электрод TFT-матрицы. Напряжение
$U$ создаёт в жидкости поле, смещающее заряженные частицы вдоль
межэлектродной оси и формирующее изображение; смена знака $U$ переключает
оптическое состояние пикселя между «белым» и «чёрным» в зависимости от того,
какие частицы оказываются у верхнего электрода.

\begin{figure}[H]
    \centering
    \begin{tikzpicture}[font=\small, scale=1.0]
        % Внешний контур микрокапсулы (круг)
        \draw[fill=blue!5, draw=blue!50!black, thick] (0, 0) circle (3);
        \node[above=0.1cm] at (0, 3) {Микрокапсула, $d_{\mathrm{caps}} \approx 40$~мкм};

        % Верхний и нижний электроды (касательные плоскости)
        \draw[ultra thick, cyan!60!black] (-3.2, 2.5) -- (3.2, 2.5);
        \node[right=0.1cm] at (3.2, 2.5) {ITO (+)};
        \draw[ultra thick, brown] (-3.2, -2.5) -- (3.2, -2.5);
        \node[right=0.1cm] at (3.2, -2.5) {TFT (--)};

        % Белые частицы (TiO2) — внутри, разбросаны
        \foreach \x/\y in {-1.5/1.5, -0.8/1.0, 0.6/1.7, 1.5/1.2, -0.5/0.5, 1.1/0.3, 0.0/-0.4, -1.7/-0.6, 1.7/-1.2, 0.3/-1.5} {
            \draw[fill=white, draw=black] (\x, \y) circle (0.18);
            \node[font=\tiny] at (\x, \y) {$+$};
        }

        % Чёрные частицы
        \foreach \x/\y in {-1.0/-1.5, 0.8/-1.7, 1.3/0.8, -1.8/0.8, 0.5/0.8, -0.2/1.4, 1.0/-0.9} {
            \draw[fill=black] (\x, \y) circle (0.18);
            \node[font=\tiny, color=white] at (\x, \y) {$-$};
        }

        % Стрелки сил
        \draw[->, very thick, red!70!black] (-2.5, 0) -- (-2.5, 1.5);
        \node[left] at (-2.5, 0.75) {$\vec{F}_E = q\vec{E}$};

        \draw[->, thick, blue!70!black] (2.5, 0.8) -- (2.5, -0.2);
        \node[right] at (2.5, 0.3) {$\vec{F}_{\!\text{Stokes}}$};

        % Поле E
        \node[font=\large] at (0, 2.1) {$\vec{E} \downarrow$};
    \end{tikzpicture}
    \caption{Сферическая микрокапсула электрофоретического дисплея с двумя
        типами заряженных пигментных частиц (белые $\mathrm{TiO}_2$ ---
        положительно заряженные, чёрные углеродные --- отрицательно
        заряженные). Под действием электрического поля $\vec{E}$ на частицу
        действуют электрическая сила $\vec{F}_E = q\vec{E}$ и противоположно
        направленная сила сопротивления среды по закону Стокса
        $\vec{F}_{\!\text{Stokes}} = -6\pi\eta R \vec{v}$.}
    \label{fig:model-capsule}
\end{figure}

Движение частицы определяется балансом электрической силы
$\vec{F}_E = q\vec{E}$ и стоксова сопротивления среды
$\vec{F}_{\mathrm{drag}} = -6\pi\eta R\, \vec{v}$ ($\eta$ — вязкость, $\vec{v}$ —
скорость). В установившемся режиме силы уравновешиваются:
$\vec{v} = q\vec{E}/(6\pi\eta R)$. Время выхода на режим задаёт постоянная
$\tau_{\mathrm{Stokes}} = m/(6\pi\eta R)$ ($m$ — масса частицы), при типовых
параметрах~--- единицы миллисекунд.

В неоднородном поле добавляется диэлектрофоретический вклад
$\vec{F}_{\mathrm{DEP}} = 2\pi r^3 \varepsilon_m\, \operatorname{Re}(K^*) \nabla|\vec{E}|^2$,
существенный для систем с пассивной адресацией без выраженного порога
переключения \cite{bert2003_physics}. В режиме настоящей работы
(последовательность постоянных по знаку импульсов в пределах одной фазы
waveform с обнулённым суммарным зарядом за обновление) им можно пренебречь;
это допущение разбирается в~\ref{sec:model-pde}, а его правомерность подтверждает
согласие упрощённой модели с измерениями (гл.~\ref{ch:experiment}).


\section{Уравнения Пуассона--Нернста--Планка}\label{sec:model-pde}

Континуальное описание концентраций двух типов частиц $c_\pm(x, t)$ строится
на законе сохранения числа частиц, конститутивном соотношении
для потока (сумма диффузионного и дрейфового слагаемых) и уравнении Пуассона,
связывающем распределение заряда с потенциалом.

\subsection*{Уравнение непрерывности}

Из закона сохранения числа частиц типа $\pm$ для произвольного объёма
$\Omega \subset \mathbb{R}^3$ имеем:
\begin{equation}\label{eq:continuity-integral}
    \frac{\partial}{\partial t} \int_{\Omega} c_\pm \, dV
    = -\oint_{\partial \Omega} \vec{J}_\pm \cdot \hat{n}\, dS,
\end{equation}
где $\vec{J}_\pm$ — поток частиц (число частиц через единицу площади в единицу
времени), $\hat{n}$ — внешняя нормаль к $\partial\Omega$. По теореме
Гаусса--Остроградского и в силу произвольности $\Omega$~--- локальная
форма:
\begin{equation}\label{eq:continuity}
    \frac{\partial c_\pm}{\partial t} + \nabla \cdot \vec{J}_\pm = 0.
\end{equation}
Уравнение~\eqref{eq:continuity} не зависит от природы потока; конкретизация
модели — в определении $\vec{J}_\pm$.

\subsection*{Конститутивное соотношение для потока}

При двух движущих силах — градиенте концентрации (диффузия) и
электрическом поле (электрофорез) — поток представим суммой:
\begin{equation}\label{eq:flux}
    \vec{J}_\pm = -D_\pm \nabla c_\pm + c_\pm \vec{v}_\pm,
\end{equation}
где первое слагаемое — закон Фика с коэффициентом $D_\pm$, второе —
конвективный перенос со средней дрейфовой скоростью $\vec{v}_\pm$.

\subsection*{Стоксова дрейфовая скорость}

Для шарообразной частицы радиуса $R_\pm$ с зарядом $q_\pm$ в жидкости вязкости
$\eta$ под действием поля $\vec{E}$ уравнение движения имеет вид:
\begin{equation}\label{eq:newton}
    m_\pm \frac{d\vec{v}_\pm}{dt}
    = q_\pm \vec{E} - 6\pi \eta R_\pm \, \vec{v}_\pm.
\end{equation}
В установившемся режиме ($d\vec{v}_\pm / dt = 0$) баланс сил даёт:
\begin{equation}\label{eq:drift-velocity}
    \vec{v}_\pm = \frac{q_\pm \vec{E}}{6\pi \eta R_\pm}
    \equiv \mu_{e,\pm}\, \vec{E},
    \qquad
    \mu_{e,\pm} \equiv \frac{|q_\pm|}{6\pi \eta R_\pm},
\end{equation}
где $\mu_{e,\pm}$ — \emph{электрофоретическая подвижность}. Знак $q_\pm$ задаёт
направление дрейфа: положительные частицы движутся вдоль поля, отрицательные —
против.

Характерное время релаксации скорости к~\eqref{eq:drift-velocity}
определяется уравнением~\eqref{eq:newton}:
\begin{equation}\label{eq:tau-stokes}
    \tau_{\mathrm{Stokes}} = \frac{m_\pm}{6\pi \eta R_\pm}.
\end{equation}
При типовых параметрах ($m \sim 10^{-17}$~кг, $R \sim 200$~нм,
$\eta \sim 10^{-3}$~Па$\cdot$с) $\tau_{\mathrm{Stokes}}$ составляет единицы
миллисекунд — значительно меньше длительностей фаз waveform
SSD1680. Поэтому далее динамика пикселя описывается в приближении
большого затухания (\emph{overdamped}; см.~\ref{sec:model-ode}).

\subsection*{Связь Эйнштейна}

В тепловом равновесии флуктуационно-диссипационная теорема связывает
коэффициент диффузии и подвижность (соотношение Эйнштейна):
\begin{equation}\label{eq:einstein}
    D_\pm = \mu_{e,\pm} \cdot \frac{k_B T}{|q_\pm|}
          = \frac{k_B T}{6\pi \eta R_\pm}.
\end{equation}
При $T \approx 295$~К, $\eta \approx 10^{-3}$~Па$\cdot$с, $R \approx 200$~нм
получаем $D \sim 10^{-12}$~м$^2$/с, что характерно для коллоидных систем.

\subsection*{Замена поля потенциалом}

Поле выражается через скалярный потенциал $\varphi(x, t)$:
\begin{equation}\label{eq:field-potential}
    \vec{E} = -\nabla \varphi.
\end{equation}
Подстановка в~\eqref{eq:flux} с учётом~\eqref{eq:drift-velocity}
и~\eqref{eq:field-potential} даёт:
\begin{equation}\label{eq:flux-final}
    \vec{J}_\pm = -D_\pm \nabla c_\pm \mp \mu_{e,\pm} c_\pm \nabla \varphi.
\end{equation}

\subsection*{Уравнения Нернста--Планка}

Подстановка~\eqref{eq:flux-final} в~\eqref{eq:continuity} при постоянных
$D_\pm$, $\mu_{e,\pm}$ в объёме капсулы (изотропная жидкость) даёт
уравнения Нернста--Планка:
\begin{equation}\label{eq:np}
    \boxed{\;\frac{\partial c_\pm}{\partial t}
    = \nabla \cdot \left(D_\pm \nabla c_\pm \pm \mu_{e,\pm} c_\pm \nabla \varphi\right).\;}
\end{equation}

\subsection*{Уравнение Пуассона}

Система~\eqref{eq:np} не замкнута: потенциал $\varphi$ зависит от распределения
заряда, определяемого концентрациями $c_\pm$. Замыкание даёт уравнение
Пуассона:
\begin{equation}\label{eq:poisson}
    -\nabla^2 \varphi = \frac{\rho_e}{\varepsilon_m \varepsilon_0},
    \qquad
    \rho_e = e\,(c_+ - c_-),
\end{equation}
где $\rho_e$ — объёмная плотность заряда, $\varepsilon_0$ — электрическая
постоянная, $\varepsilon_m$ — относительная проницаемость жидкости.

Совместная система~\eqref{eq:np} и~\eqref{eq:poisson} образует задачу
\emph{Пуассона--Нернста--Планка} (PNP) — фундаментальный объект
электрокинетики коллоидов и электролитов.

\subsection*{Граничные и начальные условия}

\textit{На электродах} ставятся условия Дирихле для потенциала и Неймана для
потоков:
\begin{align}
    \varphi(x = 0, t) &= 0,
    & \varphi(x = L, t) &= U(t), \label{eq:bc-dirichlet} \\
    \vec{J}_\pm \cdot \hat{n} \big|_{x=0,L} &= 0. \label{eq:bc-neumann}
\end{align}
Условие~\eqref{eq:bc-dirichlet} задаёт управляющее напряжение $U(t)$ —
\emph{переменную управления} в задаче гл.~\ref{ch:theorem};
условие~\eqref{eq:bc-neumann} отражает отсутствие реакций на электродах и
поглощения или генерации частиц на границах.

\textit{Начальное состояние} — однородное распределение частиц в капсуле:
\begin{equation}
    c_\pm(x, 0) = \frac{N_\pm}{V_{\mathrm{caps}}},
    \label{eq:ic}
\end{equation}
где $N_\pm$ — полное число частиц каждого типа, $V_{\mathrm{caps}}$ — объём
микрокапсулы.

\subsection*{Замечание о пренебрежении диэлектрофорезом}

Полная формулировка с диэлектрофорезом добавляла бы в правую
часть~\eqref{eq:np} член $\sim \nabla \cdot \left(c_\pm \chi \nabla |\nabla\varphi|^2\right)$,
где $\chi$ зависит от поляризуемости частицы и среды \cite{bert2003_physics}.
В режиме настоящей работы (постоянные по знаку импульсы в пределах фазы
waveform с обнулённым суммарным зарядом) им можно пренебречь. При постоянном
напряжении $|\vec{E}|$ в объёме практически постоянно (кроме узких пограничных
слоёв порядка дебаевского радиуса), так что $\nabla |\vec{E}|^2 \approx 0$;
к тому же симметрия положительных и отрицательных импульсов в сбалансированном
waveform взаимно гасит накапливаемый эффект за одно обновление. Согласованность
редуцированной модели с измерениями (гл.~\ref{ch:experiment}) подтверждает
обоснованность пренебрежения.


\section{Mean-field редукция к ОДУ одного пикселя}\label{sec:model-ode}

Прямое численное решение~\eqref{eq:np}--\eqref{eq:poisson} в задачах большой
размерности (LUT контроллера SSD1680 содержит $\sim 150$
параметров, см. гл.~\ref{ch:theorem}) практически нереализуемо. Здесь
выполняется \emph{редукция методом среднего поля} (mean-field reduction)
PNP-системы к ОДУ для центра масс распределения частиц одного пикселя.
Вычислительно она несравнимо проще исходной, но удерживает её качественное
ядро: электрофоретический дрейф, релаксацию Стокса, зарядовый баланс.

\subsection*{Допущения редукции методом среднего поля}

\textbf{Допущение~A1 (одномерная аппроксимация капсулы).} Сферическая капсула
приближается осью симметрии $x \in [0, L]$, где $L = d_{\mathrm{caps}}$ —
межэлектродное расстояние; коэффициент сферичности абсорбируется
в эффективную проницаемость поля и длину пути частицы. Допущение оправдано
осевой симметрией задачи при равномерном по горизонтали распределении частиц.

\textbf{Допущение~A2 (сосредоточенное положение, \emph{lumped position}).}
Распределение каждого типа частиц приближается дельта-функцией около центра
масс
$\bar{x}_\pm(t) = \frac{\int_0^L x\, c_\pm(x, t)\, dx}{\int_0^L c_\pm(x, t)\, dx}$,
т.\,е. отслеживание плотности $c_\pm(x, t)$ сводится к одному скаляру
$\bar{x}_\pm(t) \in [0, L]$. Допущение справедливо в режиме большого затухания,
когда диффузионное расплывание медленнее дрейфа; такой режим в EPD реализуется
при типичных напряжениях и температурах.

\textbf{Допущение~A3 (квазистатический потенциал, \emph{quasi-static}).}
Перераспределение зарядов под действием $U(t)$ идёт на временах
$\tau_{\mathrm{RC}} \sim \varepsilon_m \varepsilon_0 / \sigma_{\mathrm{eff}}$,
много меньших $\tau_{\mathrm{Stokes}}$. Поэтому при характерном времени
waveform $\sim 100$~мс потенциал $\varphi(x, t)$ устанавливается
мгновенно по текущему $c_\pm(x, t)$, что сводит~\eqref{eq:poisson}
к алгебраическому уравнению при фиксированном $t$.

\subsection*{Получаемая ОДУ-система}

Усреднением~\eqref{eq:np} с весом $x$ по $[0, L]$ при допущениях A1--A3 получаем
уравнение динамики центра масс. В режиме большого затухания (когда
$m_\pm \ddot{x}_\pm \ll 6\pi\eta R_\pm \dot{x}_\pm$ при
$t \gg \tau_{\mathrm{Stokes}}$) оно приводится к первому порядку:
\begin{equation}\label{eq:ode-position}
    \frac{d\bar{x}_\pm}{dt}
    = \pm \frac{q_\pm}{6\pi \eta R_\pm} \cdot \frac{U(t)}{L}
    = \pm \mu_{e,\pm} \cdot \frac{U(t)}{L}.
\end{equation}
В полной форме (с инерционным членом, без приближения большого затухания):
\begin{equation}\label{eq:ode-full}
    m_\pm \frac{d^2 \bar{x}_\pm}{dt^2}
    = q_\pm \frac{U(t)}{L} - 6\pi \eta R_\pm \frac{d\bar{x}_\pm}{dt}.
\end{equation}

\begin{remark}
Здесь $\bar{x}_\pm$~--- физическое положение (в метрах, $\bar{x} \in [0, L]$).
Начиная с гл.~\ref{ch:theorem} мы переходим к \emph{безразмерному}
положению $\bar{x} \mapsto \bar{x}/L \in [0, 1]$; множитель $\mu^{*}/L$
в дискретной динамике тогда приобретает размерность $1/(\text{В}\cdot\text{с})$, а
приращение положения становится безразмерным. Эта нормировка используется во
всех формулировках и доказательствах гл.~\ref{ch:theorem} и
соответствующих приложений.
\end{remark}

В практических расчётах используется~\eqref{eq:ode-position}, однако симулятор
(гл.~\ref{ch:simulation}) реализует обе формы, позволяя проверить
приближение большого затухания для конкретной комбинации параметров.

\subsection*{Преобразование положения частицы в отражательную способность}

Оптическая характеристика пикселя --- \emph{отражательная способность}
(reflectance) $\rho \in [\rho_{\min}, \rho_{\max}]$ — линейно связывается с
положением белых частиц:
\begin{equation}\label{eq:rho-x}
    \rho\bigl(\bar{x}_+\bigr) = \rho_{\min}
    + (\rho_{\max} - \rho_{\min}) \cdot
    \frac{\bar{x}_+ - x_{\mathrm{bot}}}{x_{\mathrm{top}} - x_{\mathrm{bot}}}.
\end{equation}
$\rho_{\min}$ отвечает белым частицам у нижнего электрода (видны чёрные),
$\rho_{\max}$ — у верхнего (виден белый). Типичные значения для b/w EPD:
$\rho_{\min} \approx 0{,}07$, $\rho_{\max} \approx 0{,}40$ (contrast
ratio $\approx 6$ \cite{yang2021_mechanisms}).

Линейная аппроксимация~\eqref{eq:rho-x} достаточна для оптимизации waveform с
точностью $\sim 5\%$~\cite{he2020_particle_activation, lin2024_low_power}.

\begin{figure}[H]
    \centering
    \begin{tikzpicture}[font=\small,
        every node/.style={align=center},
        level/.style={draw, very thick, rounded corners, minimum width=4.0cm, minimum height=1.4cm, fill=#1!10},
        arrow/.style={->, very thick, >=Stealth}]

        % 3 уровня по центру
        \node[level=blue]   (pde)      at (0,  4.0) {\textbf{PDE}\\континуальный};
        \node[level=green]  (ode)      at (0,  1.0) {\textbf{ОДУ}\\усреднённый};
        \node[level=orange] (discrete) at (0, -2.0) {\textbf{Дискрет.}\\по фазам waveform};

        % Стрелки между уровнями
        \draw[arrow] (pde)  -- (ode);
        \draw[arrow] (ode)  -- (discrete);

        % Подписи стрелок справа (сжали: меньше отступ, уже text width)
        \node[right=0.8cm of pde, anchor=west, text width=4.2cm, font=\footnotesize\itshape] (lbl-pde-ode)
            {\textbf{редукция методом\\среднего поля:}\\
             A1 — 1D ось симметрии\\
             A2 — сосредоточенное\\
             ~~~~~~~ положение\\
             A3 — квазистат. $\varphi$};
        \draw[->, gray, thin] (lbl-pde-ode.west) to[bend right=10] ($(pde.east)!0.5!(ode.east) + (0.2, 0)$);

        \node[right=0.8cm of ode, anchor=west, text width=4.2cm, font=\footnotesize\itshape] (lbl-ode-disc)
            {\textbf{дискретизация:}\\
             $U(t) \to \{(V_k, T_k)\}_{k=1}^{K}$\\
             $V_k \in \mathcal{U}$, $|\mathcal{U}| = 5$};
        \draw[->, gray, thin] (lbl-ode-disc.west) to[bend right=10] ($(ode.east)!0.5!(discrete.east) + (0.2, 0)$);

        % Слева — формулы (так же сжато)
        \node[left=0.8cm of pde, anchor=east, text width=4.2cm, font=\footnotesize\itshape, align=right]
            {$\partial c_\pm/\partial t = \nabla\!\cdot(\ldots)$\\
             $-\nabla^2 \varphi = \rho_e/(\varepsilon_m\varepsilon_0)$\\
             граничные условия};
        \node[left=0.8cm of ode, anchor=east, text width=4.2cm, font=\footnotesize\itshape, align=right]
            {$\dot{\bar{x}} = \mu^* \, U(t)/L$\\
             $\rho(\bar{x})$ — отражат. способн.};
        \node[left=0.8cm of discrete, anchor=east, text width=4.2cm, font=\footnotesize\itshape, align=right]
            {$\bar{x}[k{+}1] = F(\bar{x}[k], V_k, T_k)$\\
             см. главу~\ref{ch:theorem}};
    \end{tikzpicture}
    \caption{Каскадная редукция модели: континуальная PDE-задача
        Пуассона\,--\,Нернста\,--\,Планка $\to$ усреднённая система ОДУ для
        центра масс пигментных частиц одного пикселя $\to$ дискретная задача
        оптимального управления над архитектурой LUT контроллера SSD1680.
        Каждая стрелка снабжена явным набором допущений редукции (см. текст
        раздела).}
    \label{fig:model-reduction}
\end{figure}
\FloatBarrier


\section{Стоксова форма для одной частицы}\label{sec:model-stokes}

Уравнение~\eqref{eq:ode-full} для одной частицы при ступенчатом напряжении
$U(t) = U_0 = \mathrm{const}$ имеет аналитическое решение, систематически
используемое в литературе по EPD
\cite{wang2022_red_ghost, he2020_particle_activation}. Перепишем~\eqref{eq:ode-full}
в виде:
\begin{equation}\label{eq:stokes-newton}
    m \frac{dv}{dt} = qE - 6\pi\eta R\, v,
\end{equation}
где $v = d\bar{x}/dt$, $E = U_0/L$. Это линейное ОДУ первого порядка
относительно $v$; интегрирующим множителем при $v(0) = 0$ получаем:
\begin{equation}\label{eq:v-of-t}
    v(t) = v_\infty \cdot \left(1 - e^{-t/\tau_{\mathrm{Stokes}}}\right),
\end{equation}
где введены обозначения:
\begin{equation}\label{eq:vinf-tau}
    v_\infty = \frac{qU_0}{6\pi\eta R L}
    \quad \text{(terminal velocity)},
    \qquad
    \tau_{\mathrm{Stokes}} = \frac{m}{6\pi\eta R}.
\end{equation}

Соотношение~\eqref{eq:v-of-t} эквивалентно формуле
из~\cite{wang2022_red_ghost} (формула 2) и~\cite{he2020_particle_activation}
(уравнение 4). Положение частицы получаем интегрированием~\eqref{eq:v-of-t}:
\begin{equation}\label{eq:x-of-t}
    \bar{x}(t) = \bar{x}(0)
    + v_\infty \cdot \left[ t - \tau_{\mathrm{Stokes}} (1 - e^{-t/\tau_{\mathrm{Stokes}}}) \right].
\end{equation}

При $t \gg \tau_{\mathrm{Stokes}}$ скорость практически выходит на
$v_\infty$, и движение становится равномерным со скоростью
$v_\infty \propto U_0$. При $t \ll \tau_{\mathrm{Stokes}}$ движение
ускоряется по линейному закону $v(t) \approx (qU_0/mL) t$.

\subsection*{Свободные параметры модели}

В правую часть~\eqref{eq:vinf-tau} входят заряд $q$, радиус $R$, масса $m$,
вязкость $\eta$ и длина $L$, но \emph{раздельно} они не появляются: симулятор и
оптимизация используют лишь две эффективные комбинации,
\begin{align}
    \mu^*  &= \frac{v_\infty}{U_0} = \frac{q}{6\pi\eta RL}
              \quad \text{(эффективная электрофоретическая подвижность)},
              \label{eq:mu-star} \\
    \tau^* &= \tau_{\mathrm{Stokes}} = \frac{m}{6\pi\eta R}
              \quad \text{(эффективное время релаксации)}.
              \label{eq:tau-star}
\end{align}
Для пересчёта положения в отражательную способность нужны ещё $\rho_{\min}$ и
$\rho_{\max}$. Все четыре параметра $(\mu^*, \tau^*, \rho_{\min}, \rho_{\max})$
извлекаются из калибровки в~\ref{sec:model-calibration} без отдельных оценок
$q$, $\eta$, $R$, $m$.


\section{Калибровка параметров по данным стенда}\label{sec:model-calibration}

Модель содержит зависящие от панели эффективные параметры
$(\mu^*, \tau^*, \rho_{\min}, \rho_{\max}, L)$. Геометрические и динамические
($\mu^*\sim10^{-9}$~м/(В$\cdot$с)~\cite{wang2022_red_ghost},
$\tau^*\approx60$~мс~\cite{he2020_particle_activation},
$\rho_{\min}=0{,}07$, $\rho_{\max}=0{,}40$~\cite{yang2021_mechanisms},
$L\approx40$~мкм) приняты по литературе для EPD сходной геометрии и служат для
качественного анализа структуры оптимума (гл.~\ref{ch:theorem}).

\emph{Энергетические} параметры калибруются \emph{операционно} по измерениям
стенда (ESP32 + Waveshare~2.13''~V4 + INA219, $R_{\mathrm{shunt}}=0{,}1$~Ом).
Энергия обновления определяется током драйверов по всей матрице, поэтому
моделируется на уровне \emph{панели}~--- как функции $E(N)$ и $\tau(N)$ числа
кадров waveform. Вывод, коэффициенты ($R^2\approx1$) и
проверка~--- в гл.~\ref{ch:experiment}, \S\ref{sec:calib}. Предварительные
одно-фазные прогоны подтверждают порядок величин: энергия обновления — единицы
мДж (1{,}5--1{,}7~мДж на одно-фазный тест-импульс), что согласуется с
данными Lin~и~соавторов~\cite{lin2024_low_power} после пересчёта на площадь
панели.

\subsection*{Границы применимости и поправка на бистабильность}

ОДУ-описание~\eqref{eq:ode-position} обосновано в overdamped-режиме и линейной
аппроксимации $\rho(\bar x)$ (точной до контраста $\sim5\%$). Эксперимент
(гл.~\ref{ch:experiment}) выявляет важную поправку: реальная панель
\emph{бистабильна}~--- итоговое состояние пикселя задаёт последняя фаза драйва, а
не интеграл $\int V\,dt$. Поэтому линейный вывод <<строгий зарядовый баланс
$\sum V T=0 \Rightarrow \Delta\bar x=0$>> на бистабильной панели не выполняется:
зарядовый баланс служит \emph{долговечности} (отсутствию DC-деградации), а не
определяет смещение. В задаче оптимизации (гл.~\ref{ch:theorem}) это
формализуется параметром~$\varepsilon$ зарядового баланса, образующим ось Парето
<<энергия--долговечность>> (гл.~\ref{ch:experiment}). Структурные результаты
глав~\ref{ch:theorem}--\ref{ch:algorithm} (релейность, оценка числа фаз) от этой
поправки не зависят и подтверждаются экспериментально.
